% !PRINTED 693
\weberpage{693}
%<<CONT
and $\varrho$ could not be divisible by $r\varepsilon$. In the same way one concludes
that each of the prime factors $\pi_1, \pi_2, \ldots$ divides $r\varepsilon$, but none
of them more than once; hence, apart from unit factors,

\weq{15}{r\varepsilon = \varkappa.}

Thus each of the functions $\varepsilon_1, \varepsilon_2, \ldots, \varepsilon_n$ satisfies condition a).
Conversely, every function $w$ that satisfies this condition is also
contained in the form:

\weq{16}{w = x_1\varepsilon_1 + x_2\varepsilon_2 + \cdots + x_n\varepsilon_n,}

in which the $x_1, x_2, \ldots, x_n$ are integral rational functions of $z$.

For $rw$ is an integral function divisible by $\varkappa$, and
consequently, by (9), (16) and § 184, 7.:
\[
S(rw\omega_i) = rx_i
\]
is an integral function of $z$ divisible by $r$; consequently

\weq{17}{x_i = S(w\omega_i).}

In order, then, to obtain the differential quotients of the first
kind, one has to select, among the functions (16), those that
satisfy condition b).

To this end we take for (2) the normal basis considered in § 193,
\[
\lambda_1, \lambda_2, \ldots, \lambda_n
\]
and denote the basis complementary to it by
\[
\mu_1, \mu_2, \ldots, \mu_n,
\]
which is defined by

\weq{18}{S(\lambda_r\mu_s) = (r, s)}

If we then set

\weq{19}{w = y_1\mu_1 + y_2\mu_2 + \cdots + y_n\mu_n,}

then, if $r_i$ is the exponent of $\lambda_i$, it follows by (17) that
\[
\frac{y_i}{z^{r_i-1}} = S\!\left(wz\frac{\lambda_i}{z^{r_i}}\right)\!.
\]

By the definition of $r_i$, $\lambda_i:z^{r_i}$ is finite for $z=\infty$, and
$wz$ vanishes on account of b). Consequently $y_i:z^{r_i-1}$ must vanish
for $z=\infty$, that is, $y_i$ can be at most of degree $(r_i-2)$
and therefore contains at most $r_i-1$ constant coefficients.

Since $\mu_1=1$, $r_1=0$ and the degree of $y_1$ cannot be
% [sic] printed $\mu_1=1$, but the fact used is $\lambda_1=1$ ($r_1=0$ is the exponent of $\lambda_1$, and $y_1=S(w\lambda_1)=S(w)$ gives $S(w)=0$); with $\lambda_1=1$, (18) gives $S(\mu_1)=1\neq S(1)=n$, so $\mu_1\neq1$
negative, $y_1$ must be identically zero, and it follows that
\[
S(w) = 0.
\]
